\documentclass[10pt]{beamer}
\setbeamertemplate{navigation symbols}{}
\usecolortheme{beaver}
\setbeamertemplate{blocks}[rounded=true, shadow=true]
\setbeamertemplate{footline}{%
  \hfill\usebeamercolor[fg]{page number in head/foot}%
  \footnotesize\insertframenumber{}/\inserttotalframenumber\hspace{3mm}\vskip2mm
}

\usepackage[utf8]{inputenc}
\usepackage[T2A]{fontenc}
\usepackage[english,russian]{babel}
\usepackage{amsmath,amssymb,amsfonts}
\usepackage{booktabs}
\usepackage{graphicx}
\usepackage{tikz}
\usetikzlibrary{arrows.meta, positioning, calc}

\DeclareMathOperator{\rank}{rank}
\DeclareMathOperator{\trace}{tr}
\DeclareMathOperator{\KL}{KL}
\DeclareMathOperator*{\argmin}{arg\,min}
\DeclareMathOperator*{\argmax}{arg\,max}
\newcommand{\R}{\mathbb{R}}
\newcommand{\E}{\mathbb{E}}
\newcommand{\Hsym}{\mathcal{H}_{\mathrm{sym}}}
\newcommand{\Hasym}{\mathcal{H}_{\mathrm{asym}}}

\title{Снижение размерности латентного пространства в задачах управления}
\author{Астахов А.\,М.}
\institute{МФТИ}
\date{2026}

\begin{document}


\begin{frame}[plain]
\begin{center}
{\Large\textbf{Снижение размерности латентного пространства\\[1pt]
в задачах управления}}\\[12mm]
{\large Астахов А.\,М.}\\[2mm]
группа Б05-251л\\[5mm]
\textit{Научный руководитель:} д.\,ф.-м.\,н.\ Стрижов В.\,В.\\[8mm]
\small
Московский физико-технический институт\\
Физтех-школа прикладной математики и информатики\\
Кафедра интеллектуальных систем\\[4mm]
2026
\end{center}
\end{frame}


\begin{frame}{Снижение размерности латентного пространства в задачах управления}
\framesubtitle{Цель, задача, решение}
\footnotesize

\textbf{Цель.}\quad Построить компактную модель прогноза для распределённой
динамической системы, наблюдаемой через сеть из $C$ датчиков, --- модель,
пригодную для использования в контуре управления в реальном времени.

\medskip
\textbf{Задача.}\quad Временной ряд каждого датчика содержит избыточные
задержечные координаты; их число нужно сократить до компактного латентного
представления. Одновременно требуется восстановить оператор связности между
датчиками --- матрицу их попарных направленных влияний --- без
предположения о её симметричности.

\medskip
\textbf{Решение.}\quad Латентное представление пониженной размерности,
согласованное с геометрией исходной траектории, и параметризованный оператор
связности, для которого симметричность не является структурным
ограничением модели.
\end{frame}


\begin{frame}{Снижение размерности латентного пространства в задачах управления}
\framesubtitle{Формальная постановка задачи}
\footnotesize

Фазовое пространство --- $\mathcal M = L^2(\Omega)$, $\Omega \subset
\R^n$, $n \in \{1,2,3\}$. Предполагается существование компактного
$\varphi$-инвариантного множества $\mathcal A \subset \mathcal M$ конечной
бокс-размерности $d_{\mathcal A} < \infty$, к которому сходятся все
траектории; его элементы $u \in \mathcal A$ --- скрытые состояния системы,
наблюдаемые через вектор показаний $s(t) = (h_1(u(t)), \dots, h_C(u(t)))$.

Скрытое состояние $u(t)$ напрямую не наблюдается. Теорема Такенса о
вложении по задержкам гарантирует, что при $Cm \ge 2 d_{\mathcal A}+1$
отображение
\[
\Phi(t) = \bigl(s_i(t-j\tau)\bigr)_{i=1,\dots,C;\,j=0,\dots,m-1} \in \R^{C \times m}
\]
является вложением $\mathcal A$ в евклидово пространство, то есть
гомеоморфизмом $\Psi: \mathcal A \to \Psi(\mathcal A)$. Поэтому оператор
связности можно параметризовать не от недоступного $u$, а от латентного
представления $z = z(\Phi(u))$, построенного по наблюдаемым задержечным
координатам; его эволюция приближённо линейна и описывается искомым
оператором связности.

\begin{block}{Задачи}
По траектории $\mathcal D = \{s(t_n)\}_{n=0}^{N}$ требуется построить
латентное представление $z(t) \in \R^{C \times L}$, $L<m$, согласованное с
теоремой Такенса; построить параметризованное отображение $\hat A_\theta: z
\mapsto \R^{C \times C}$; обучить модель прогноза $s(t+k\Delta t)$ на
горизонте $H$ итеративной динамикой в латентном пространстве.
\end{block}
\end{frame}


\begin{frame}{Допущения и структурные ограничения}
\footnotesize

Любая матрица $A \in \R^{C \times C}$ единственным образом раскладывается в
ортогональную сумму симметричной и антисимметричной частей:
\[
\R^{C \times C} = \Hsym \oplus \Hasym, \qquad
A = A_{\mathrm{sym}} + A_{\mathrm{asym}}, \qquad
\|A\|_F^2 = \|A_{\mathrm{sym}}\|_F^2 + \|A_{\mathrm{asym}}\|_F^2.
\]
Корреляция Пирсона $\hat A = \gamma \widehat\Sigma$, ядерные графы
$\kappa(\|\xi_i-\xi_j\|)$ и self-attention с $W^Q=W^K$ принадлежат $\Hsym$ по
построению.

\begin{block}{Неустранимая ошибка симметричного класса}
Для любой $A$ и любой $M \in \Hsym$ выполнено $\|A-M\|_F^2 \ge
\|A_{\mathrm{asym}}\|_F^2$, с равенством при $M=A_{\mathrm{sym}}$. Если
истинный оператор $A(u)$ имеет ненулевую антисимметричную часть, любая
симметричная модель несёт среднюю ошибку $\E\,\|A_{\mathrm{asym}}(u)\|_F^2$,
не зависящую ни от объёма данных, ни от числа параметров: модель
$\hat A_\theta(z)$ не должна быть ограничена подпространством $\Hsym$.
\end{block}

Дополнительно принимается допущение о низком ранге: операторы связности
физических систем, как правило, эффективно низкоранговы, $\rank(A) \le L$.
Параметризация с этим ограничением сокращает число параметров оператора с
$C^2$ до $2CL$ и повышает устойчивость оценки при ограниченном объёме данных.
\end{frame}

\begin{frame}{Архитектура MoE-GA}
\centering
\resizebox{\textwidth}{!}{%
\begin{tikzpicture}[
    node distance=6mm,
    box/.style={draw, rounded corners=2pt, align=center, font=\footnotesize,
                 inner sep=4pt, text width=3.0cm, minimum height=1.3cm},
    line/.style={-{Latex[length=2mm]}, thick},
]
\node[box] (tk) {Такенс\\$\Phi(t)\in\R^{C\times m}$};
\node[box, right=of tk] (ae) {Автоэнкодер\\$z(t)\in\R^{C\times L}$};
\node[box, right=of ae] (temp) {Временной\\трансформер};

\node[box, below=of temp] (moe) {Смесь экспертов\\+ Гумбель-роутер};
\node[box, left=of moe] (ga) {GA-шаг\\(residual)};
\node[box, left=of ga] (dec) {Декодер\\$\hat s(t+k\Delta t)$};

\draw[line] (tk) -- (ae);
\draw[line] (ae) -- (temp);
\draw[line] (temp) -- (moe);
\draw[line] (moe) -- (ga);
\draw[line] (ga) -- (dec);
\end{tikzpicture}%
}

\bigskip
\footnotesize
\raggedright
Задержечные векторы Такенса сжимаются автоэнкодером в латент
$z(t)\in\R^{C\times L}$ размерности $L=\mathrm{round}(2m/3)$; временной
трансформер обрабатывает латентную траекторию по каждому датчику отдельно,
накапливая временной контекст.

\medskip
Полученный латент поступает в смесь $K$ независимо параметризованных
экспертов: маршрутизатор Гумбеля жёстко выбирает один эксперт и строит
асимметричный оператор связности $\hat A(z)$. Шаг graph attention обновляет
латент как residual: $z \leftarrow z + \alpha\,\hat A(z)V(z)$, в полной
версии \emph{MoE-GA Full} дополнительно учитывается реакционный член
$f_\theta(z)$. Линейный декодер восстанавливает прогноз показаний.
\end{frame}


\begin{frame}{Асимметричное внимание и маршрутизация Гумбеля}
\footnotesize

Для каждой моды $k=1,\dots,K$ заведены две независимые проекции $W^Q_k,
W^K_k \in \R^{L\times d}$, $d=L$. Модовая матрица --- нормированная
билинейная форма
\[
A_k(z) = \frac{(zW^Q_k)\,(zW^K_k)^\top}{\sqrt d} \in \R^{C\times C}.
\]
Поскольку $W^Q_k$ и $W^K_k$ независимы, $A_k(z)$ в общем случае
несимметрична, а её ранг не превышает $d=L$: низкоранговое допущение с
предыдущего слайда выполняется здесь автоматически по построению.

Маршрутизатор $R_\psi$ берёт усреднённый по датчикам латент $\bar z =
\frac1C\sum_i z_i$ и выдаёт логиты $h(z)=R_\psi(\bar z)\in\R^K$. При обучении
используется Gumbel--Softmax с убывающей температурой; на выводе --- жёсткий
выбор $\hat k(z)=\argmax_k h_k(z)$, активирующий ровно одного эксперта.
Итоговый оператор $\hat A(z) = \mathrm{SN}\bigl(A_{\hat k(z)}(z)\bigr)$
нормируется по спектральному радиусу до целевого значения $r=1{,}5$.

\begin{block}{Симметричный класс как частный случай}
При $W^Q_k=W^K_k$ матрица $A_k(z)$ становится симметричной матрицей Грама,
$\hat A(z)\in\Hsym$. При несвязанных проекциях и $d\ge C$ образ отображения
$z\mapsto A_k(z)$ покрывает весь $\R^{C\times C}$, так что
$\mathcal F_\Theta^{\mathrm{sym}} \subsetneq \mathcal F_\Theta$: построенный
класс операторов строго шире симметричного.
\end{block}
\end{frame}


\begin{frame}{Совместная оптимизационная задача}
\footnotesize
Энкодер детерминирован: $z(t)=\mu_\phi(\Phi(t))$, без сэмплирования.
Реконструкционный член
\[
\mathcal L_{\mathrm{rec}} = \frac{1}{|\mathcal T_{\mathrm{tr}}|\,Cm}
\sum_{t\in\mathcal T_{\mathrm{tr}}} \|\Phi(t) - D_\phi(z(t))\|_2^2
\]
обеспечивает информативность латента; компактность задаётся
ограничением размерности $L<m$.

Пространственный блок обучается прогнозными лоссами: горизонтный
$\mathcal L_{\mathrm{pred}}$ (rollout на $H$ шагов) и внутриоконный
$\mathcal L_{\mathrm{ind}}$, дающий $W-1$ точек на окно. Регуляризатор
маршрутизации $\mathcal L_{\mathrm{reg}} = \KL(\mathrm{Uniform}\{1,\dots,K\}
\,\|\,\bar w(B)) \ge 0$ предотвращает коллапс в один эксперт;
регуляризатор фазы $\mathcal L_{\mathrm{phase}} =
\mathrm{BCE}(\sigma(\hat r_\chi(\bar z(t))), r(t))$ обучает голову
$\hat r_\chi$ предсказывать бинарный режим $r(t)$, делая $z(t)$
режимо-различающим.
\[
\mathcal L(\theta) = \mathcal L_{\mathrm{pred}} + \lambda_1 \mathcal
L_{\mathrm{ind}} + \lambda_2 \mathcal L_{\mathrm{rec}} + \lambda_3
\mathcal L_{\mathrm{reg}} + \lambda_4 \mathcal L_{\mathrm{phase}}, \qquad
\theta^\ast \in \argmin_{\theta\in\Theta} \mathcal L(\theta).
\]

\begin{block}{Согласованность задачи}
Все слагаемые обучаются совместно AdamW с косинусным расписанием
learning rate. Член $\mathcal L_{\mathrm{pred}}$ заставляет $\hat
A_\theta(z)$ точно описывать направленную динамику; $\mathcal
L_{\mathrm{rec}}$ при $L<m$ обеспечивает компактное латентное
представление $CL \ll Cm$; $\mathcal L_{\mathrm{reg}}$ и $\mathcal
L_{\mathrm{phase}}$ (последний --- только на синтетике, на ЭЭГ
$\equiv0$) поддерживают равномерную и режимо-осознанную загрузку $K$
экспертов. Снижение размерности и обучение оператора связности
происходят согласованно.
\end{block}
\end{frame}


\begin{frame}{Эксперимент 1: кольцо с дрейфующим ядром (drift-ring)}
\footnotesize
\begin{columns}[T]
\begin{column}{0.55\textwidth}
Одномерное кольцо из $C=16$ датчиков; эволюция задана циркулянтным
оператором с гауссовым пиком, центр которого переключается каждые
$500$ шагов между $+v_0$ и $-v_0$ --- бинарное чередование режимов
переноса вправо/влево. Истинная асимметрия в каждом режиме
$\|T_\pm-T_\pm^\top\|_F/\|T_\pm\|_F \approx 1{,}41$, усреднённый по
режимам оператор симметричен. Шум $\sigma=0{,}3$, $T=3000$,
$L=3$, $K=6$ (вручную, перекрывает два физических режима).

\medskip
\centering
\setlength{\tabcolsep}{3pt}
\begin{tabular}{lccc}
\toprule
Модель & MSE & $A_{\mathrm{asym}}$ & Улучш. \\
\midrule
Корреляция   & $0{,}229_{\pm0{,}035}$ & $\approx 0$ & --- \\
MoE-GA       & $0{,}211_{\pm0{,}023}$ & $0{,}467$   & $7{,}9\%$ \\
MoE-GA Full  & $0{,}204_{\pm0{,}024}$ & $0{,}255$   & $11{,}1\%$ \\
\bottomrule
\end{tabular}
\end{column}
\begin{column}{0.45\textwidth}
\centering
\includegraphics[width=\linewidth]{figures/drift/fig_summary_bar.png}

\smallskip
\raggedright
Симметричный baseline даёт $A_{\mathrm{asym}}\approx 0$ по построению ---
в точности неустранимая ошибка из утверждения о пределе симметричного
класса. MoE-GA Full снижает MSE на $11{,}1\%$, восстанавливая
асимметрию $0{,}255$; обе модели MoE-GA устойчиво опережают корреляцию
на всех шагах горизонта $h=1,\dots,4$.
\end{column}
\end{columns}
\end{frame}


\begin{frame}{Восстановленные матрицы связности: drift-ring}
\centering
\includegraphics[height=3.9cm,width=\textwidth,keepaspectratio]{figures/drift/fig1_regimes.png}

\smallskip
\footnotesize
\raggedright
Средние матрицы $\hat A$ по двум режимам ($v>0$ и $v<0$) для трёх
моделей и истинные операторы $T_+$, $T_-$. \emph{Корреляция}
симметрична по построению ($a=0$) и одинаково коррелирует с обоими
режимами ($\rho=0{,}33$): направление переноса не различается, что и есть
неустранимая ошибка симметричного класса. \emph{MoE-GA} восстанавливает
асимметрию $a=1{,}58$, близкую к истинной $1{,}41$, и сильнее
коррелирует со «своим» режимом ($\rho(\hat A_+,T_+)=0{,}51$,
$\rho(\hat A_-,T_-)=0{,}37$); маршрутизатор активирует преимущественно
две моды из шести, отвечающие двум направлениям переноса.
\emph{MoE-GA Full} имеет малую асимметрию матрицы ($a=0{,}09$), так как
часть направленной динамики перехватывает реакционный член, но знаки
корреляций с $T_+$ и $T_-$ противоположны ($0{,}45$ и $-0{,}43$) ---
инверсия направления при смене режима обнаруживается корректно.
\end{frame}


\begin{frame}{Эксперимент 2: ЭЭГ моторного воображения}
\footnotesize
\begin{columns}[T]
\begin{column}{0.55\textwidth}
Один испытуемый PhysioNet EEGBCI, эксперименты $4,8,12$ (motor imagery
левой/правой руки), $C=64$ электрода, частота $160$~Гц, $T=30000$
отсчётов после полосовой фильтрации $1$--$40$~Гц и стандартизации.
Гиперпараметры выбираются автоматически по данным: $\tau=9$, $m=5$,
$L=\mathrm{round}(5\times2/3)=3$; критерий Марченко--Пастура даёт
$K_{\mathrm{MP}}=4$, итоговое $K=4$. Горизонт прогноза $H=8$.

\medskip
\centering
\setlength{\tabcolsep}{3pt}
\begin{tabular}{lccc}
\toprule
Модель & MSE & $A_{\mathrm{asym}}$ & Улучш. \\
\midrule
Корреляция   & $0{,}472_{\pm0{,}007}$ & $\approx 0$ & --- \\
MoE-GA       & $0{,}450_{\pm0{,}006}$ & $0{,}016$   & $4{,}7\%$ \\
MoE-GA Full  & $0{,}440_{\pm0{,}007}$ & $0{,}085$   & $6{,}8\%$ \\
\bottomrule
\end{tabular}
\end{column}
\begin{column}{0.45\textwidth}
\centering
\includegraphics[width=\linewidth]{figures/eeg/fig_horizon.png}

\smallskip
\raggedright
Прирост скромнее, чем на синтетике: направленная составляющая в ЭЭГ
моторного воображения выражена слабее, основная часть прогностически
значимой информации сосредоточена в симметричной части. По горизонту
прогноза MoE-GA Full опережает остальные варианты на длинных шагах
($h\ge3$), при близких ошибках на первом шаге.
\end{column}
\end{columns}
\end{frame}

\begin{frame}{Восстановленные матрицы связности: ЭЭГ}
\centering
\includegraphics[height=3.6cm,width=\textwidth,keepaspectratio]{figures/eeg/fig_eeg_A_mean.png}

\bigskip
\footnotesize
\raggedright
Истинный оператор связи для ЭЭГ недоступен, поэтому анализ носит
структурный характер. Усреднённые по всем сэмплам матрицы $\hat A$
($C=64$): \emph{Corr} (слева, $a=0$), \emph{MoE-GA} (центр,
$a=0{,}023$), \emph{MoE-GA Full} (справа, $a=0{,}040$). Во всех трёх
видна блочная структура, соответствующая анатомическим группам
электродов --- моторная кора (Cz, C3, C4), лобные и затылочные
кластеры, что согласуется с нейрофизиологическими ожиданиями для
задачи motor imagery. Маршрутизатор MoE-GA Full почти равномерно
распределяет $K=4$ режима ($\bar w \approx [0{,}31,\,0{,}38,\,0{,}14,\,
0{,}17]$), что согласуется с критерием $K_{\mathrm{MP}}=4$; доминирующий
эксперт кодирует низкочастотную $\delta$-компоненту, остальные ---
более локальные связи.
\end{frame}

\begin{frame}{Выносится на защиту}
\footnotesize
\raggedright

Симметричные модели связности --- корреляция, ядерные графы,
tied-attention --- несут неустранимую ошибку $\E\,\|A_{\mathrm{asym}}(u)\|_F^2$
при наличии направленных взаимодействий между датчиками.

\bigskip
Латентное представление $z(t)\in\R^{C\times L}$, $L<m$, построенное по
теореме Такенса, достаточно для прогноза динамики распределённой системы.

\bigskip
Параметризованный оператор связности $\hat A_\theta(z)$ строится вне
подпространства $\Hsym$ и сохраняет ранг $\rank(\hat A_\theta) \le L$.

\bigskip
Архитектура MoE-GA снижает ошибку прогноза на $11{,}1\%$ (drift-ring,
$C=16$) и на $6{,}8\%$ (ЭЭГ, $C=64$) относительно симметричной
корреляционной модели.

\bigskip
Восстановленные операторы связности согласуются со структурой исходной
системы: режимами переноса в drift-ring и анатомическими группами
электродов в ЭЭГ.
\end{frame}

\end{document}
